\documentclass[11pt]{article}
\usepackage[utf8]{inputenc}  
\usepackage[english]{babel}    
\usepackage{hyperref}        
\usepackage{booktabs}		 
\usepackage{a4wide}          
\pagestyle{empty}            
\usepackage{graphicx}        
\usepackage{grffile}         
\usepackage{caption}         
\usepackage{float}
\usepackage{fancyhdr}
\usepackage{titling}
\setlength{\droptitle}{-2cm} 
\posttitle{\end{center}\vspace{-0.5cm}} 
\postauthor{\end{center}\vspace{-0.5cm}} 
\predate{\begin{center}}
\postdate{\end{center}\vspace{-0.8cm}} 

\title{Digital Electronic Systems Design Project Report}
\author{Federico Ferro, Simone Guan, Xiao Chen Wang, Cheng Marco Zhang}
\date{Academic Year 2025/2026} 

\begin{document}

\maketitle

\fancypagestyle{plain}{
    \fancyhf{}
    \lfoot{\footnotesize Politecnico di Milano}
    \rfoot{\footnotesize Spring 2026: Digital Electronic Systems Design}
    \renewcommand{\headrulewidth}{0pt}
    \renewcommand{\footrulewidth}{0pt}
}

\section*{Timing Analysis}

The implemented design demonstrates the timing analysis yields compliant results for both setup and hold constraints at the target operating frequency of 150.89286~MHz.

Regarding the \textbf{setup time}, the most critical path (Path~21) is located within the 
\texttt{reverb\_0} sub-block. While the design requirement demands a 150.89286~MHz clock, this implementation achieves a positive slack of 0.299~ns. This safety margin allows the circuit to potentially operate at a maximum frequency of approximately 158~MHz, derived as $f_{\max} = 1\,/\,(6.627\,\mathrm{ns} - 0.299\,\mathrm{ns}) \approx 158.0\,\mathrm{MHz}$.

Concerning the \textbf{hold time}, the most critical path (Path~31) is located within the 
\texttt{axis\_fifo\_0} sub-block. The required hold time constraint is 0~ns, and the implementation guarantees a safe positive hold slack of 0.024~ns, ensuring that no race 
conditions occur.

The detailed parameters resulting from the implementation are summarized in Table~\ref{tab:setup_report} for setup time and in Table~\ref{tab:hold_report} for hold time below, related to the above mentioned most critical paths.

\begin{table}[H]
    \centering 
    \begin{minipage}[b]{0.48\textwidth}
        \flushleft
        \captionsetup{type=table, singlelinecheck=false, justification=raggedright, skip=2pt}
        \caption{Timing Summary: Setup Time (Path 21)}
        \label{tab:setup_report}
        \small
        \begin{tabular}{lrc}
            \toprule
            \textbf{Parameter} & \textbf{Value} & \textbf{Unit} \\
            \midrule
            Slack             & 0.299 & ns \\
            Requirement       & 6.627 & ns \\
            Data Path Delay   & 2.378 & ns \\
            \quad \textit{Logic Delay} & 0.742 & ns \\
            \quad \textit{Route Delay} & 1.636 & ns \\
            Clock Path Skew   & 0.055 & ns \\
            Clock Uncertainty & 0.079 & ns \\
            \bottomrule
        \end{tabular}
    \end{minipage}%
    \hfill 
    \begin{minipage}[b]{0.48\textwidth}
        \flushleft
        \captionsetup{type=table, singlelinecheck=false, justification=raggedright, skip=2pt}
        \caption{Timing Summary: Hold Time (Path 31)}
        \label{tab:hold_report}
        \small
        \begin{tabular}{lrc}
            \toprule
            \textbf{Parameter} & \textbf{Value} & \textbf{Unit} \\
            \midrule
            Slack             & 0.024 & ns \\
            Requirement       & 0.000 & ns \\
            Data Path Delay   & 0.347 & ns \\
            \quad \textit{Logic Delay} & 0.141 & ns \\
            \quad \textit{Route Delay} & 0.206 & ns \\
            Clock Path Skew   & 0.013 & ns \\
            \bottomrule
        \end{tabular}
    \end{minipage}
    
\end{table}

\section*{Power Analysis}

At the operating frequency of 150.89286~MHz, the on-chip power consumption of the implemented design totals 0.334~W, broken down into a \textbf{dynamic power} component of 0.261~W (78\%) and a \textbf{static power} (device leakage) component of 0.074~W (22\%). The dominant contributor to dynamic power is the MMCM clock manager at 0.126~W (48\% of dynamic), followed by BRAM at 0.078~W (30\%) and I/O at 0.017~W (6\%). The junction temperature of the package is estimated at 26.7~\textdegree C, with a thermal margin of 58.3~\textdegree C and an effective thermal resistance $\theta_{JA}$ of 5.0~\textdegree C/W, indicating comfortable thermal headroom under the given operating conditions. The confidence level of the power estimate is reported as Low, as no switching activity data from simulation was provided, and the analysis relies on vectorless estimation. The detailed power figures are reported in Table~\ref{tab:main_power} and Table~\ref{tab:power_dynamic_breakdown}.

\begin{table}[H]
    \centering 
    \begin{minipage}[b]{0.46\textwidth}
        \flushleft
        \captionsetup{type=table, singlelinecheck=false, justification=raggedright, skip=2pt}
        \caption{Power Analysis Summary.}
        \label{tab:main_power}
        \small
        \begin{tabular}{lrc}
            \toprule
            \textbf{Parameter} & \textbf{Value} & \textbf{Unit} \\
            \midrule
            Total On-Chip Power            & 0.334 & W   \\
            Design Power Budget            & N/A   & --  \\
            Power Budget Margin            & N/A   & --  \\
            Junction Temperature           & 26.7  & °C  \\
            Thermal Margin                 & 58.3  & °C  \\
            Effective $\theta_{JA}$        & 5.0   & °C/W\\
            Power to Off-Chip Devices      & 0     & W   \\
            Confidence Level               & Low   & --  \\
            \bottomrule
        \end{tabular}
    \end{minipage}%
    \hfill 
    \begin{minipage}[b]{0.50\textwidth}
        \flushleft
        \captionsetup{type=table, singlelinecheck=false, justification=raggedright, skip=2pt}
        \caption{On-Chip Power Breakdown.}
        \label{tab:power_dynamic_breakdown}
        \small
        \begin{tabular}{lrcc}
            \toprule
            \textbf{Component} & \textbf{Power} & \textbf{Unit} & \textbf{Ratio (\%)} \\
            \midrule
            \textbf{Dynamic}   & \textbf{0.261} & \textbf{W} & \textbf{78\%} \\
            \quad \textit{MMCM}               & 0.126 & W & 48\% \\
            \quad \textit{BRAM}               & 0.078 & W & 30\% \\
            \quad \textit{I/O}                & 0.017 & W &  6\% \\
            \quad \textit{Clocks}             & 0.018 & W &  7\% \\
            \quad \textit{Signals}            & 0.012 & W &  5\% \\
            \quad \textit{Logic}              & 0.009 & W &  3\% \\
            \quad \textit{DSP}                & 0.002 & W &  1\% \\
            \midrule
            \textbf{Device Static} & \textbf{0.074} & \textbf{W} & \textbf{22\%} \\
            \bottomrule
        \end{tabular}
    \end{minipage}
    
\end{table}

\end{document}
